% Include this table with: % Include this table with: % Include this table with: % Include this table with: \input{table_comparison}
%
% Training regime comparison table for FluidWorld on Moving MNIST.
% Final results: Pixel (30ep), JEPA (30ep), Random baseline.

\begin{table}[t]
    \centering
    \caption{\textbf{Training regime comparison on Moving MNIST.} Three training objectives applied to the same \fluidworld{} architecture ($\sim$801K parameters), evaluated after 30 epochs (Random baseline uses untrained weights). $^\dagger$The ``JEPA-style'' column refers to \fluidworld{} trained with a latent prediction objective inspired by LeCun's JEPA framework~\cite{lecun2022path}: the pixel decoder is removed from the loss and replaced by a target-encoder (EMA) + VICReg objective. This is \emph{not} Meta's I-JEPA or V-JEPA architecture. Latent cosine similarity is the primary discriminating metric: it measures whether the model's internal rollout tracks the true future state in representation space.}
    \label{tab:training_regime_comparison}
    \begin{tabular}{lccc}
        \toprule
        \textbf{Metric} & \textbf{Pixel (30 ep)} & \textbf{JEPA-style$^\dagger$ (30 ep)} & \textbf{Random} \\
        \midrule
        \multicolumn{4}{l}{\textit{Linear probes on frozen BeliefField features}} \\
        \quad Position $R^2$ (spatial)$^a$    & 0.9766          & 0.9349          & \textbf{0.9522} \\
        \quad Velocity $R^2$ (after PDE)$^b$  & \textbf{0.7730} & 0.2876          & 0.5580          \\
        \quad Velocity $R^2$ (MLP probe)$^{b2}$ & 0.6328         & \textbf{0.6028} & 0.3145          \\
        \quad Direction Acc.\ (8-class)$^c$   & \textbf{38.4\%} & 22.1\%          & 30.8\%          \\
        \midrule
        \multicolumn{4}{l}{\textit{Latent-space rollout quality (cosine similarity $\uparrow$)}$^d$} \\
        \quad Step 1                          & 0.054           & \textbf{0.833}  & 0.060           \\
        \quad Step 10                         & 0.143           & \textbf{0.827}  & 0.069           \\
        \quad Step 19                         & $-$0.002        & \textbf{0.828}  & 0.047           \\
        \bottomrule
    \end{tabular}

    \vspace{0.5em}
    \footnotesize
    $^a$ Position $R^2$ is non-discriminating: even random weights achieve 0.97, since spatial position is recoverable from any convolutional feature map via linear readout. \\
    $^b$ The JEPA-style model's low \emph{linear} velocity $R^2$ reflects non-linear encoding of velocity in representation space, not failure. Cosine similarity of 0.828 at step~19 and the MLP probe result (footnote $b_2$) confirm this. \\
    $^{b2}$ With a 2-layer MLP probe, the JEPA-style model's velocity $R^2$ rises from 0.288 to \textbf{0.603} ($\Delta = +0.315$), while Pixel drops from 0.773 to 0.633 ($\Delta = -0.14$) and Random drops from 0.558 to 0.315 ($\Delta = -0.24$). The JEPA-style model is the only one whose velocity $R^2$ increases under a non-linear probe, consistent with distributed, non-linearly entangled representations. Linear probes underestimate its encoding of dynamics. \\
    $^c$ Chance level is 12.5\% for 8 direction classes. Pixel's advantage here reflects linearly decodable velocity, not superior dynamics modeling. \\
    $^d$ Latent cosine similarity measures whether the model's autonomous rollout tracks the true future. The JEPA-style model maintains $\geq 0.827$ through step~19; Pixel is near zero and Random fluctuates around chance. It produces no pixel-space predictions because no decoder was trained (by design, it operates in representation space). \\
    $^e$ Edge/Freq auxiliary losses (Sobel + FFT) were tested as an ablation. At 30 epochs, Edge/Freq degraded early-step SSIM from 0.778 to 0.013 relative to Laplacian-only. See ablation analysis for details.
\end{table}

\begin{table}[t]
    \centering
    \caption{\textbf{FluidWorld in the context of the JEPA research program.} Y. LeCun's 2022 cognitive architecture blueprint\protect\footnotemark[2] identified several open challenges. Meta FAIR's V-JEPA~2 tackles large-scale video understanding; \fluidworld{} explores complementary directions (rollout stability, self-repair, physics-grounded dynamics) at a much smaller scale. The two approaches are not in competition but address different parts of the same research agenda.}
    \label{tab:fluidworld_vs_jepa}
    \begin{tabular}{lp{4.8cm}p{4.8cm}}
        \toprule
        \textbf{Dimension} & \textbf{Meta FAIR (V-JEPA 2)\protect\footnotemark[1]} & \textbf{FluidWorld} \\
        \midrule
        Model size
            & $\sim$1.3B parameters
            & $\sim$801K parameters \\
        \addlinespace
        Latent rollout stability (step 19)
            & Error accumulation (acknowledged bottleneck)
            & Cosine 0.827 (stable) \\
        \addlinespace
        Autopoietic self-repair
            & Outside current scope
            & Recovery 66.8\%, $p < 10^{-49}$ \\
        \addlinespace
        Perturbation robustness
            & Outside current scope
            & 50\% corruption recovered in 3--7 steps \\
        \addlinespace
        Physics-informed structure
            & Data-driven (ViT)
            & Reaction-diffusion PDE \\
        \addlinespace
        Persistent memory
            & Fixed context window
            & Titans + DeltaNet \\
        \addlinespace
        Persistent state (BeliefField)
            & Proposed (Y. LeCun 2022\protect\footnotemark[2]), not implemented
            & Implemented ($16 \times 16 \times 128$) \\
        \addlinespace
        Memory complexity
            & $O(N^2)$
            & $O(N)$ \\
        \addlinespace
        Training hardware
            & 16+ A100 GPUs
            & 1 RTX 4070 Ti \\
        \addlinespace
        Edge of chaos dynamics
            & Outside current scope
            & $\lambda = 0.0033$ (supercritical) \\
        \bottomrule
    \end{tabular}

    \vspace{0.5em}
    \footnotetext[1]{Bardes et al., ``V-JEPA 2: Self-Supervised Video Models Enable Understanding, Prediction, and Planning,'' arXiv:2506.09985, 2025.}
    \footnotetext[2]{LeCun, ``A Path Towards Autonomous Machine Intelligence,'' OpenReview, 2022.}
\end{table}

%
% Training regime comparison table for FluidWorld on Moving MNIST.
% Final results: Pixel (30ep), JEPA (30ep), Random baseline.

\begin{table}[t]
    \centering
    \caption{\textbf{Training regime comparison on Moving MNIST.} Three training objectives applied to the same \fluidworld{} architecture ($\sim$801K parameters), evaluated after 30 epochs (Random baseline uses untrained weights). $^\dagger$The ``JEPA-style'' column refers to \fluidworld{} trained with a latent prediction objective inspired by LeCun's JEPA framework~\cite{lecun2022path}: the pixel decoder is removed from the loss and replaced by a target-encoder (EMA) + VICReg objective. This is \emph{not} Meta's I-JEPA or V-JEPA architecture. Latent cosine similarity is the primary discriminating metric: it measures whether the model's internal rollout tracks the true future state in representation space.}
    \label{tab:training_regime_comparison}
    \begin{tabular}{lccc}
        \toprule
        \textbf{Metric} & \textbf{Pixel (30 ep)} & \textbf{JEPA-style$^\dagger$ (30 ep)} & \textbf{Random} \\
        \midrule
        \multicolumn{4}{l}{\textit{Linear probes on frozen BeliefField features}} \\
        \quad Position $R^2$ (spatial)$^a$    & 0.9766          & 0.9349          & \textbf{0.9522} \\
        \quad Velocity $R^2$ (after PDE)$^b$  & \textbf{0.7730} & 0.2876          & 0.5580          \\
        \quad Velocity $R^2$ (MLP probe)$^{b2}$ & 0.6328         & \textbf{0.6028} & 0.3145          \\
        \quad Direction Acc.\ (8-class)$^c$   & \textbf{38.4\%} & 22.1\%          & 30.8\%          \\
        \midrule
        \multicolumn{4}{l}{\textit{Latent-space rollout quality (cosine similarity $\uparrow$)}$^d$} \\
        \quad Step 1                          & 0.054           & \textbf{0.833}  & 0.060           \\
        \quad Step 10                         & 0.143           & \textbf{0.827}  & 0.069           \\
        \quad Step 19                         & $-$0.002        & \textbf{0.828}  & 0.047           \\
        \bottomrule
    \end{tabular}

    \vspace{0.5em}
    \footnotesize
    $^a$ Position $R^2$ is non-discriminating: even random weights achieve 0.97, since spatial position is recoverable from any convolutional feature map via linear readout. \\
    $^b$ The JEPA-style model's low \emph{linear} velocity $R^2$ reflects non-linear encoding of velocity in representation space, not failure. Cosine similarity of 0.828 at step~19 and the MLP probe result (footnote $b_2$) confirm this. \\
    $^{b2}$ With a 2-layer MLP probe, the JEPA-style model's velocity $R^2$ rises from 0.288 to \textbf{0.603} ($\Delta = +0.315$), while Pixel drops from 0.773 to 0.633 ($\Delta = -0.14$) and Random drops from 0.558 to 0.315 ($\Delta = -0.24$). The JEPA-style model is the only one whose velocity $R^2$ increases under a non-linear probe, consistent with distributed, non-linearly entangled representations. Linear probes underestimate its encoding of dynamics. \\
    $^c$ Chance level is 12.5\% for 8 direction classes. Pixel's advantage here reflects linearly decodable velocity, not superior dynamics modeling. \\
    $^d$ Latent cosine similarity measures whether the model's autonomous rollout tracks the true future. The JEPA-style model maintains $\geq 0.827$ through step~19; Pixel is near zero and Random fluctuates around chance. It produces no pixel-space predictions because no decoder was trained (by design, it operates in representation space). \\
    $^e$ Edge/Freq auxiliary losses (Sobel + FFT) were tested as an ablation. At 30 epochs, Edge/Freq degraded early-step SSIM from 0.778 to 0.013 relative to Laplacian-only. See ablation analysis for details.
\end{table}

\begin{table}[t]
    \centering
    \caption{\textbf{FluidWorld in the context of the JEPA research program.} Y. LeCun's 2022 cognitive architecture blueprint\protect\footnotemark[2] identified several open challenges. Meta FAIR's V-JEPA~2 tackles large-scale video understanding; \fluidworld{} explores complementary directions (rollout stability, self-repair, physics-grounded dynamics) at a much smaller scale. The two approaches are not in competition but address different parts of the same research agenda.}
    \label{tab:fluidworld_vs_jepa}
    \begin{tabular}{lp{4.8cm}p{4.8cm}}
        \toprule
        \textbf{Dimension} & \textbf{Meta FAIR (V-JEPA 2)\protect\footnotemark[1]} & \textbf{FluidWorld} \\
        \midrule
        Model size
            & $\sim$1.3B parameters
            & $\sim$801K parameters \\
        \addlinespace
        Latent rollout stability (step 19)
            & Error accumulation (acknowledged bottleneck)
            & Cosine 0.827 (stable) \\
        \addlinespace
        Autopoietic self-repair
            & Outside current scope
            & Recovery 66.8\%, $p < 10^{-49}$ \\
        \addlinespace
        Perturbation robustness
            & Outside current scope
            & 50\% corruption recovered in 3--7 steps \\
        \addlinespace
        Physics-informed structure
            & Data-driven (ViT)
            & Reaction-diffusion PDE \\
        \addlinespace
        Persistent memory
            & Fixed context window
            & Titans + DeltaNet \\
        \addlinespace
        Persistent state (BeliefField)
            & Proposed (Y. LeCun 2022\protect\footnotemark[2]), not implemented
            & Implemented ($16 \times 16 \times 128$) \\
        \addlinespace
        Memory complexity
            & $O(N^2)$
            & $O(N)$ \\
        \addlinespace
        Training hardware
            & 16+ A100 GPUs
            & 1 RTX 4070 Ti \\
        \addlinespace
        Edge of chaos dynamics
            & Outside current scope
            & $\lambda = 0.0033$ (supercritical) \\
        \bottomrule
    \end{tabular}

    \vspace{0.5em}
    \footnotetext[1]{Bardes et al., ``V-JEPA 2: Self-Supervised Video Models Enable Understanding, Prediction, and Planning,'' arXiv:2506.09985, 2025.}
    \footnotetext[2]{LeCun, ``A Path Towards Autonomous Machine Intelligence,'' OpenReview, 2022.}
\end{table}

%
% Training regime comparison table for FluidWorld on Moving MNIST.
% Final results: Pixel (30ep), JEPA (30ep), Random baseline.

\begin{table}[t]
    \centering
    \caption{\textbf{Training regime comparison on Moving MNIST.} Three training objectives applied to the same \fluidworld{} architecture ($\sim$801K parameters), evaluated after 30 epochs (Random baseline uses untrained weights). $^\dagger$The ``JEPA-style'' column refers to \fluidworld{} trained with a latent prediction objective inspired by LeCun's JEPA framework~\cite{lecun2022path}: the pixel decoder is removed from the loss and replaced by a target-encoder (EMA) + VICReg objective. This is \emph{not} Meta's I-JEPA or V-JEPA architecture. Latent cosine similarity is the primary discriminating metric: it measures whether the model's internal rollout tracks the true future state in representation space.}
    \label{tab:training_regime_comparison}
    \begin{tabular}{lccc}
        \toprule
        \textbf{Metric} & \textbf{Pixel (30 ep)} & \textbf{JEPA-style$^\dagger$ (30 ep)} & \textbf{Random} \\
        \midrule
        \multicolumn{4}{l}{\textit{Linear probes on frozen BeliefField features}} \\
        \quad Position $R^2$ (spatial)$^a$    & 0.9766          & 0.9349          & \textbf{0.9522} \\
        \quad Velocity $R^2$ (after PDE)$^b$  & \textbf{0.7730} & 0.2876          & 0.5580          \\
        \quad Velocity $R^2$ (MLP probe)$^{b2}$ & 0.6328         & \textbf{0.6028} & 0.3145          \\
        \quad Direction Acc.\ (8-class)$^c$   & \textbf{38.4\%} & 22.1\%          & 30.8\%          \\
        \midrule
        \multicolumn{4}{l}{\textit{Latent-space rollout quality (cosine similarity $\uparrow$)}$^d$} \\
        \quad Step 1                          & 0.054           & \textbf{0.833}  & 0.060           \\
        \quad Step 10                         & 0.143           & \textbf{0.827}  & 0.069           \\
        \quad Step 19                         & $-$0.002        & \textbf{0.828}  & 0.047           \\
        \bottomrule
    \end{tabular}

    \vspace{0.5em}
    \footnotesize
    $^a$ Position $R^2$ is non-discriminating: even random weights achieve 0.97, since spatial position is recoverable from any convolutional feature map via linear readout. \\
    $^b$ The JEPA-style model's low \emph{linear} velocity $R^2$ reflects non-linear encoding of velocity in representation space, not failure. Cosine similarity of 0.828 at step~19 and the MLP probe result (footnote $b_2$) confirm this. \\
    $^{b2}$ With a 2-layer MLP probe, the JEPA-style model's velocity $R^2$ rises from 0.288 to \textbf{0.603} ($\Delta = +0.315$), while Pixel drops from 0.773 to 0.633 ($\Delta = -0.14$) and Random drops from 0.558 to 0.315 ($\Delta = -0.24$). The JEPA-style model is the only one whose velocity $R^2$ increases under a non-linear probe, consistent with distributed, non-linearly entangled representations. Linear probes underestimate its encoding of dynamics. \\
    $^c$ Chance level is 12.5\% for 8 direction classes. Pixel's advantage here reflects linearly decodable velocity, not superior dynamics modeling. \\
    $^d$ Latent cosine similarity measures whether the model's autonomous rollout tracks the true future. The JEPA-style model maintains $\geq 0.827$ through step~19; Pixel is near zero and Random fluctuates around chance. It produces no pixel-space predictions because no decoder was trained (by design, it operates in representation space). \\
    $^e$ Edge/Freq auxiliary losses (Sobel + FFT) were tested as an ablation. At 30 epochs, Edge/Freq degraded early-step SSIM from 0.778 to 0.013 relative to Laplacian-only. See ablation analysis for details.
\end{table}

\begin{table}[t]
    \centering
    \caption{\textbf{FluidWorld in the context of the JEPA research program.} Y. LeCun's 2022 cognitive architecture blueprint\protect\footnotemark[2] identified several open challenges. Meta FAIR's V-JEPA~2 tackles large-scale video understanding; \fluidworld{} explores complementary directions (rollout stability, self-repair, physics-grounded dynamics) at a much smaller scale. The two approaches are not in competition but address different parts of the same research agenda.}
    \label{tab:fluidworld_vs_jepa}
    \begin{tabular}{lp{4.8cm}p{4.8cm}}
        \toprule
        \textbf{Dimension} & \textbf{Meta FAIR (V-JEPA 2)\protect\footnotemark[1]} & \textbf{FluidWorld} \\
        \midrule
        Model size
            & $\sim$1.3B parameters
            & $\sim$801K parameters \\
        \addlinespace
        Latent rollout stability (step 19)
            & Error accumulation (acknowledged bottleneck)
            & Cosine 0.827 (stable) \\
        \addlinespace
        Autopoietic self-repair
            & Outside current scope
            & Recovery 66.8\%, $p < 10^{-49}$ \\
        \addlinespace
        Perturbation robustness
            & Outside current scope
            & 50\% corruption recovered in 3--7 steps \\
        \addlinespace
        Physics-informed structure
            & Data-driven (ViT)
            & Reaction-diffusion PDE \\
        \addlinespace
        Persistent memory
            & Fixed context window
            & Titans + DeltaNet \\
        \addlinespace
        Persistent state (BeliefField)
            & Proposed (Y. LeCun 2022\protect\footnotemark[2]), not implemented
            & Implemented ($16 \times 16 \times 128$) \\
        \addlinespace
        Memory complexity
            & $O(N^2)$
            & $O(N)$ \\
        \addlinespace
        Training hardware
            & 16+ A100 GPUs
            & 1 RTX 4070 Ti \\
        \addlinespace
        Edge of chaos dynamics
            & Outside current scope
            & $\lambda = 0.0033$ (supercritical) \\
        \bottomrule
    \end{tabular}

    \vspace{0.5em}
    \footnotetext[1]{Bardes et al., ``V-JEPA 2: Self-Supervised Video Models Enable Understanding, Prediction, and Planning,'' arXiv:2506.09985, 2025.}
    \footnotetext[2]{LeCun, ``A Path Towards Autonomous Machine Intelligence,'' OpenReview, 2022.}
\end{table}

%
% Training regime comparison table for FluidWorld on Moving MNIST.
% Final results: Pixel (30ep), JEPA (30ep), Random baseline.

\begin{table}[t]
    \centering
    \caption{\textbf{Training regime comparison on Moving MNIST.} Three training objectives applied to the same \fluidworld{} architecture ($\sim$801K parameters), evaluated after 30 epochs (Random baseline uses untrained weights). $^\dagger$The ``JEPA-style'' column refers to \fluidworld{} trained with a latent prediction objective inspired by LeCun's JEPA framework~\cite{lecun2022path}: the pixel decoder is removed from the loss and replaced by a target-encoder (EMA) + VICReg objective. This is \emph{not} Meta's I-JEPA or V-JEPA architecture. Latent cosine similarity is the primary discriminating metric: it measures whether the model's internal rollout tracks the true future state in representation space.}
    \label{tab:training_regime_comparison}
    \begin{tabular}{lccc}
        \toprule
        \textbf{Metric} & \textbf{Pixel (30 ep)} & \textbf{JEPA-style$^\dagger$ (30 ep)} & \textbf{Random} \\
        \midrule
        \multicolumn{4}{l}{\textit{Linear probes on frozen BeliefField features}} \\
        \quad Position $R^2$ (spatial)$^a$    & 0.9766          & 0.9349          & \textbf{0.9522} \\
        \quad Velocity $R^2$ (after PDE)$^b$  & \textbf{0.7730} & 0.2876          & 0.5580          \\
        \quad Velocity $R^2$ (MLP probe)$^{b2}$ & 0.6328         & \textbf{0.6028} & 0.3145          \\
        \quad Direction Acc.\ (8-class)$^c$   & \textbf{38.4\%} & 22.1\%          & 30.8\%          \\
        \midrule
        \multicolumn{4}{l}{\textit{Latent-space rollout quality (cosine similarity $\uparrow$)}$^d$} \\
        \quad Step 1                          & 0.054           & \textbf{0.833}  & 0.060           \\
        \quad Step 10                         & 0.143           & \textbf{0.827}  & 0.069           \\
        \quad Step 19                         & $-$0.002        & \textbf{0.828}  & 0.047           \\
        \bottomrule
    \end{tabular}

    \vspace{0.5em}
    \footnotesize
    $^a$ Position $R^2$ is non-discriminating: even random weights achieve 0.97, since spatial position is recoverable from any convolutional feature map via linear readout. \\
    $^b$ The JEPA-style model's low \emph{linear} velocity $R^2$ reflects non-linear encoding of velocity in representation space, not failure. Cosine similarity of 0.828 at step~19 and the MLP probe result (footnote $b_2$) confirm this. \\
    $^{b2}$ With a 2-layer MLP probe, the JEPA-style model's velocity $R^2$ rises from 0.288 to \textbf{0.603} ($\Delta = +0.315$), while Pixel drops from 0.773 to 0.633 ($\Delta = -0.14$) and Random drops from 0.558 to 0.315 ($\Delta = -0.24$). The JEPA-style model is the only one whose velocity $R^2$ increases under a non-linear probe, consistent with distributed, non-linearly entangled representations. Linear probes underestimate its encoding of dynamics. \\
    $^c$ Chance level is 12.5\% for 8 direction classes. Pixel's advantage here reflects linearly decodable velocity, not superior dynamics modeling. \\
    $^d$ Latent cosine similarity measures whether the model's autonomous rollout tracks the true future. The JEPA-style model maintains $\geq 0.827$ through step~19; Pixel is near zero and Random fluctuates around chance. It produces no pixel-space predictions because no decoder was trained (by design, it operates in representation space). \\
    $^e$ Edge/Freq auxiliary losses (Sobel + FFT) were tested as an ablation. At 30 epochs, Edge/Freq degraded early-step SSIM from 0.778 to 0.013 relative to Laplacian-only. See ablation analysis for details.
\end{table}

\begin{table}[t]
    \centering
    \caption{\textbf{FluidWorld in the context of the JEPA research program.} Y. LeCun's 2022 cognitive architecture blueprint\protect\footnotemark[2] identified several open challenges. Meta FAIR's V-JEPA~2 tackles large-scale video understanding; \fluidworld{} explores complementary directions (rollout stability, self-repair, physics-grounded dynamics) at a much smaller scale. The two approaches are not in competition but address different parts of the same research agenda.}
    \label{tab:fluidworld_vs_jepa}
    \begin{tabular}{lp{4.8cm}p{4.8cm}}
        \toprule
        \textbf{Dimension} & \textbf{Meta FAIR (V-JEPA 2)\protect\footnotemark[1]} & \textbf{FluidWorld} \\
        \midrule
        Model size
            & $\sim$1.3B parameters
            & $\sim$801K parameters \\
        \addlinespace
        Latent rollout stability (step 19)
            & Error accumulation (acknowledged bottleneck)
            & Cosine 0.827 (stable) \\
        \addlinespace
        Autopoietic self-repair
            & Outside current scope
            & Recovery 66.8\%, $p < 10^{-49}$ \\
        \addlinespace
        Perturbation robustness
            & Outside current scope
            & 50\% corruption recovered in 3--7 steps \\
        \addlinespace
        Physics-informed structure
            & Data-driven (ViT)
            & Reaction-diffusion PDE \\
        \addlinespace
        Persistent memory
            & Fixed context window
            & Titans + DeltaNet \\
        \addlinespace
        Persistent state (BeliefField)
            & Proposed (Y. LeCun 2022\protect\footnotemark[2]), not implemented
            & Implemented ($16 \times 16 \times 128$) \\
        \addlinespace
        Memory complexity
            & $O(N^2)$
            & $O(N)$ \\
        \addlinespace
        Training hardware
            & 16+ A100 GPUs
            & 1 RTX 4070 Ti \\
        \addlinespace
        Edge of chaos dynamics
            & Outside current scope
            & $\lambda = 0.0033$ (supercritical) \\
        \bottomrule
    \end{tabular}

    \vspace{0.5em}
    \footnotetext[1]{Bardes et al., ``V-JEPA 2: Self-Supervised Video Models Enable Understanding, Prediction, and Planning,'' arXiv:2506.09985, 2025.}
    \footnotetext[2]{LeCun, ``A Path Towards Autonomous Machine Intelligence,'' OpenReview, 2022.}
\end{table}
